\subsection{Problemstellung \& Motivation}
\label{sec:problemstellung_motivation}

\begin{tcolorbox}[boxrule=0pt,sharp corners,enhanced,borderline west={2pt}{0pt}{blue}]
    \textbf{Für manche eine optionale Lesehilfe, für andere der einzige Zugang zu gesellschaftlicher Teilhabe.}
\end{tcolorbox}

Beim Thema Barrierefreiheit denken viele an physische Barrieren wie Rampen, Aufzüge oder automatische Türen. Durch die Digitalisierung und die Verlagerung von Dienstleistungen und Informationen ins Internet ist es damit jedoch nicht mehr getan. Barrierefreie Gestaltung betrifft zunehmend digitale Räume und insbesondere die Sprache, in der Informationen bereitgestellt werden.

Menschen mit kognitiven Einschränkungen, Lernschwierigkeiten, geringer Lese- und Schreibkompetenz oder Deutsch als Zweitsprache werden durch komplexe Texte, bürokratische Fachsprache und verschachtelte Satzstrukturen vom Zugang zu grundlegenden Informationen und Dienstleistungen ausgeschlossen. Das Konzept der \textit{Leichten Sprache} (LS) bietet hierfür ein klar reglementiertes Regelwerk, um Texte durch Vereinfachung des Vokabulars und der Syntax verständlich zu machen. 

Mit dem Inkrafttreten des Barrierefreiheitsstärkungsgesetzes (BFSG) in Deutschland sowie europäischen Richtlinien (wie der Web-Accessibility-Directive) wächst der rechtliche Druck auf öffentliche und private Akteure, Informationen auch in Leichter Sprache anzubieten. Die manuelle Übersetzung von Alltagssprache (AS) in Leichte Sprache (LS) durch menschliche Übersetzer ist jedoch extrem zeitintensiv, kostspielig und durch einen Mangel an Fachkräften limitiert.

Hieraus ergibt sich ein dringender Forschungs- und Entwicklungsbedarf im Bereich der maschinellen Textvereinfachung. Bisherige neuronale Übersetzungsmodelle scheitern jedoch häufig an den spezifischen syntaktischen und strukturellen Besonderheiten der Leichten Sprache oder leiden unter dem Mangel an qualitativ hochwertigen, parallelen Trainingsdaten. Es bedarf daher systematischer Ansätze zur automatisierten Erstellung paralleler Korpora, robuster Metriken zur Komplexitätsbewertung und optimierter Trainingsstrategien.

\subsection{Zielsetzung \& Forschungsfragen}
\label{sec:zielsetzung_forschungsfragen}

Ziel dieser Arbeit ist die Entwicklung, Optimierung und empirische Evaluierung eines methodischen Frameworks zur automatisierten Übersetzung deutscher Alltagssprache in Leichte Sprache. Dies umfasst die gesamte Pipeline: von der Akquise und Bereinigung paralleler Web-Korpora über die kontinuierliche, referenzfreie Modellierung von Textsimplizität bis hin zur dateneffizienten Modelladaption (LoRA) und dem Präferenz-Alignment (DPO).

Im Zentrum dieser Arbeit steht die folgende übergreifende Leitfrage:

\begin{tcolorbox}[boxrule=0pt,sharp corners,enhanced,borderline west={2pt}{0pt}{blue}]
    \textbf{\large Wie lässt sich die automatische Übersetzung in Leichte Sprache (LS) durch die Kombination aus rauschbereinigten Korpora, kontinuierlichen Simplizitätsmetriken und zielgerichteten Alignment-Strategien (DPO) systematisch optimieren und validieren?}
\end{tcolorbox}

Zur systematischen Beantwortung dieser Leitfrage werden sechs forschungsleitende Kernfragen (FF1 bis FF6) formuliert, die den logischen Bogen über die drei methodischen Hauptsäulen der Arbeit spannen:

\begin{enumerate}
    \item \textbf{Datenbasis \& Korpusqualität (Themenblock 1):}
    \begin{itemize}
        \item \textbf{FF1 (Quellenheterogenität \& Web-Rauschen):} \textit{Inwieweit existiert im deutschsprachigen Web eine hinreichende und qualitativ konsistente Datenbasis für Leichte Sprache, und welche Vorverarbeitungsschritte sind erforderlich, um domänenspezifisches Web-Rauschen (Boilerplate, Teaser, maschinelle Vorübersetzungen) zu eliminieren?}
        \item \textbf{FF2 (Artikel- vs. Satzebene \& Informationsverlust):} \textit{Welche Vor- und Nachteile bietet eine artikelbasierte Modellierung gegenüber satzweiser Segmentierung für das Übersetzen in Leichte Sprache, und wie lässt sich der inhärente Konflikt zwischen syntaktischer Simplifikation und semantischem Informationsverlust (Faktenerhalt) quantifizieren?}
    \end{itemize}
    
    \item \textbf{Metrik \& Kontinuierliche Simplizitäts-Modellierung (Themenblock 2):}
    \begin{itemize}
        \item \textbf{FF3 (Kontinuierliche Regression ohne Zwischenkorpora):} \textit{Wie kann ein kontinuierliches Regressionsmodell zur Messung des Vereinfachungsgrads trainiert werden, wenn in realen Korpora ausschließlich binäre Extremzustände (Standardsprache vs. Leichte Sprache) ohne natürliche Zwischenstufen vorliegen?}
    \end{itemize}

    \item \textbf{Übersetzung, Alignment \& Expertenvalidierung (Themenblock 3 \& Diskussion):}
    \begin{itemize}
        \item \textbf{FF4 (Dateneffizienz \& PEFT/LoRA):} \textit{In welchem Maße können vortrainierte Sequenz-zu-Sequenz-Modelle bei extrem begrenzter Datenmenge (~800--1.000 Artikelpaare) die komplexen syntaktischen und lexikalischen Regeln der Leichten Sprache erlernen, ohne in Memorization und Overfitting zu verfallen?}
        \item \textbf{FF5 (Reward-Guided Alignment via DPO):} \textit{Inwiefern kann das Signal eines kontinuierlich trainierten Simplizitäts-Regressors genutzt werden, um ein vortrainiertes Übersetzungsmodell über Direct Preference Optimization (DPO) gezielt auf die Eigenheiten Leichter Sprache auszurichten, ohne die semantische Faktentreue zu kompromittieren?}
        \item \textbf{FF6 (Expertenvalidierung \& Ökologische Validität):} \textit{Wie valide korrelieren automatisierte Bewertungsmetriken mit dem qualitativen Urteil von Fachkräften für Leichte Sprache (Lebenshilfe Kiel), und in welchem Maße sind Modellübersetzungen von professionellen Humanübersetzungen unterscheidbar?}
    \end{itemize}
\end{enumerate}

\subsection{Aufbau der Arbeit}
\label{sec:aufbau_der_arbeit}

Die vorliegende Arbeit folgt einer logischen Struktur, die sich an den drei Säulen des entwickelten Übersetzungssystems orientiert:

\begin{itemize}
    \item \textbf{Kapitel 2 (Theoretischer Hintergrund)} legt das theoretische und linguistische Fundament. Es definiert das Konzept der Leichten Sprache und grenzt dieses formal sowie rechtlich von der Einfachen Sprache ab.
    
    \item \textbf{Kapitel 3 (Stand der Forschung)} stellt den aktuellen wissenschaftlichen Stand strukturiert entlang der drei Säulen Datenbasis, Komplexitätsmetriken und maschinelle Übersetzungsverfahren dar und leitet die bestehenden Forschungslücken ab.
    
    \item \textbf{Kapitel 4 (Themenblock 1: Datenbasis \& Korpus-Erstellung)} beschreibt die erste Stufe der Pipeline. Im Fokus steht das Crawling paralleler Webtexte, die domänenspezifische Datenbereinigung und die qualitative Filterung mittels semantischer Embeddings.
    
    \item \textbf{Kapitel 5 (Themenblock 2: Metrik \& Bewerten von Sprachkomplexität)} widmet sich der zweiten Stufe: der Entwicklung von Klassifikationsmodellen (BiLSTM vs. SBERT) und dem Entwurf einer \textit{Continuous MixUp Regression} zur stufenlosen Komplexitätsbewertung.
    
    \item \textbf{Kapitel 6 (Themenblock 3: Modellierung der Übersetzung \& Reward-Guided Fine-Tuning)} stellt das eigentliche Übersetzungsmodell vor. Hierbei werden die Implementierung des Supervised Fine-Tunings (SFT) sowie das Reward-Guided Fine-Tuning (DPO) zur Regeleinhaltung evaluiert.
    
    \item \textbf{Kapitel 7 (Diskussion \& Gesamtevaluation)} führt die Ergebnisse der drei Stufen zusammen, beantwortet systematisch die aufgeworfenen Forschungsfragen und beleuchtet kritisch die Systemrestriktionen sowie Fragen zur Faktentreue.
    
    \item \textbf{Kapitel 8 (Fazit \& Ausblick)} fasst die Kernergebnisse der Arbeit zusammen und gibt einen Ausblick auf zukünftige Forschungsarbeiten und praktische Anwendungsfelder.
\end{itemize}
